\documentclass{article}

% NeurIPS style file
\usepackage{neurips_2026}

% Packages
\usepackage[utf8]{inputenc}
\usepackage[T1]{fontenc}
\usepackage{hyperref}
\usepackage{url}
\usepackage{booktabs}
\usepackage{amsfonts}
\usepackage{amsmath}
\usepackage{amssymb}
\usepackage{graphicx}
\usepackage{float}
\usepackage{xcolor}
\usepackage{enumitem}
\usepackage{array}
\usepackage{tabularx}

% Project title and authors
\title{CSC413 Project Proposal: [Your Project Title]}

\author{
  Member One \\
  University of Toronto
  \And
  Member Two \\
  University of Toronto
  \And
  Member Three \\
  University of Toronto
}

\begin{document}

\maketitle

% ============================================================
\begin{abstract}
% Target: approximately half a page.
% Describe:
% 1. The problem and motivation.
% 2. The task your model will perform.
% 3. The dataset and its important characteristics.
% 4. The proposed neural network architecture.
% 5. How you will evaluate the model.
% 6. Why the results could be useful or meaningful.

[Write your abstract here. Give the reader enough information about
the project, data, task, and model to understand the work and determine
whether they have the expertise to review it.]

\end{abstract}

% ============================================================
\section{Introduction}

% Explain WHAT, WHY, and SO WHAT.
% - What problem are you solving?
% - Why is the problem important?
% - Who might benefit from a solution?
% - Why is deep learning a reasonable approach?
% - What is your general approach?
% - What are your main goals and expected contributions?

\subsection{Problem and Motivation}

[Introduce the problem and explain its real-world or research
significance.]

\subsection{Project Goals}

[Describe the specific task, intended model behavior, and measurable
objectives.]

\subsection{Proposed Approach}

[Briefly introduce the planned model, data, and evaluation strategy.
Explain why a neural network is appropriate for the task.]

% ============================================================
\section{Background and Related Work}

% Provide enough domain context for CSC413 students and TAs.
% Discuss at least 1--2 works directly related to your task.
% Cite sources for domain concepts, datasets, and referenced models.

\subsection{Background}

[Explain relevant domain concepts, terminology, and task-specific
challenges that readers may not already know.]

\subsection{Related Work}

[Discuss at least one or two prior papers that solve similar problems.
For each paper, describe its task, approach, key findings, and
relationship to your project.]

\subsection{How Our Work Differs}

[Explain what you will adapt, compare, investigate, or contribute
beyond the existing work.]

% Example citation:
% Prior work has explored related approaches \citep{author2024}.

% ============================================================
\section{Data}

\subsection{Dataset Description}

[Describe the dataset source, size, features, labels, and data types.
Discuss class balance, missing values, noise, and other relevant
characteristics. Explain why the dataset is suitable for your task.]

\subsection{Data Preprocessing}

[Describe reproducible cleaning and formatting steps, including
filtering, normalization, tokenization, resizing, encoding, or
augmentation as appropriate.]

\subsection{Training, Validation, and Test Splits}

[Explain how the data will be split, including split proportions
and how data leakage will be prevented.]

\subsection{Exploratory Data Analysis}

[Describe important patterns, distributions, examples, or challenges
discovered during data exploration. Add plots or tables if useful.]

% ============================================================
\section{Model Architecture}

\subsection{Model Overview}

[Describe the type of neural network you plan to use and why it fits
the task and the structure of the data.]

\subsection{Architecture Components}

[Explain the key layers or components, their roles, and how information
flows through the model. Include relevant dimensions or hyperparameters
where known.]

\subsection{Input and Output}

[Specify the input representation, output format, prediction target,
and any activation functions or loss functions you plan to use.]

\subsection{Training Strategy}

[Describe the planned optimizer, learning rate, batch size, number of
epochs, regularization, and evaluation metrics. Mark undecided choices
as preliminary plans.]

% ============================================================
\section{Model Architecture Figure}

% Upload your figure to Overleaf, for example as architecture.png.
% Replace the filename below with your actual figure filename.

\begin{figure}[H]
    \centering
    % Uncomment once your figure has been uploaded.
    % \includegraphics[width=0.9\linewidth]{architecture.png}

    \fbox{
        \parbox[c][2.0in][c]{0.85\linewidth}{
            \centering
            \textbf{Architecture Figure Placeholder}\\[0.5em]
            Input Data $\rightarrow$ Preprocessing
            $\rightarrow$ Neural Network
            $\rightarrow$ Prediction
            $\rightarrow$ Evaluation
        }
    }

    \caption{Overview of the proposed model architecture.
    Replace this placeholder with a diagram showing the actual
    data flow and important model components.}
    \label{fig:architecture}
\end{figure}

[Explain the key design choices illustrated in Figure~\ref{fig:architecture}.]

% ============================================================
\section{Ethical Considerations}

\subsection{Data Collection and Privacy}

[Discuss dataset consent, privacy, sensitive information, licensing,
representation, and possible collection biases.]

\subsection{Bias and Fairness}

[Consider whether the data or model could disadvantage particular
groups or systematically produce unequal errors.]

\subsection{Potential Benefits and Harms}

[Explain beneficial applications, misuse scenarios, limitations,
and the consequences of incorrect predictions.]

\subsection{Mitigation Strategies}

[Describe practical steps to reduce risks, such as data filtering,
privacy protection, subgroup evaluation, transparency, or human review.]

% ============================================================
\section{Work Division}

% Replace the example rows with your actual team members and plans.
% Include responsibilities, weekly time estimates, deliverables,
% dependencies, and potential scheduling conflicts.

\begin{table}[H]
    \centering
    \small
    \caption{Division of responsibilities and planned time commitment.}
    \label{tab:workdivision}
    \begin{tabularx}{\linewidth}{
        @{}p{0.16\linewidth}
        X
        p{0.16\linewidth}
        p{0.18\linewidth}@{}
    }
        \toprule
        \textbf{Member} &
        \textbf{Responsibilities and Deliverables} &
        \textbf{Weeks} &
        \textbf{Hours/Week} \\
        \midrule
        Member One &
        Literature review, dataset investigation, and data preprocessing. &
        1--3 &
        [X] \\
        \addlinespace
        Member Two &
        Model implementation, architecture design, and training pipeline. &
        2--5 &
        [X] \\
        \addlinespace
        Member Three &
        Evaluation, baseline comparisons, and result visualization. &
        3--5 &
        [X] \\
        \addlinespace
        All members &
        Paper writing, integration, debugging, and final review. &
        1--6 &
        [X] \\
        \bottomrule
    \end{tabularx}
\end{table}

\subsection{Weekly Plan}

\begin{itemize}[leftmargin=*]
    \item \textbf{Week 1:} Finalize the task, research related work,
    and inspect the dataset.
    \item \textbf{Week 2:} Complete preprocessing and establish a baseline.
    \item \textbf{Week 3:} Implement the proposed model and begin training.
    \item \textbf{Week 4:} Tune hyperparameters and run experiments.
    \item \textbf{Week 5:} Analyze results, perform comparisons, and
    address limitations.
    \item \textbf{Week 6:} Finalize the report, figures, and reproducibility
    checks.
\end{itemize}

\subsection{Coordination and Dependencies}

[Explain which tasks depend on others, how code and results will be
integrated, how often the team will meet, and how conflicts or delays
will be handled.]

% ============================================================
% References
% Create a file named references.bib and add your sources there.

\bibliographystyle{plainnat}
\bibliography{references}

\end{document}
```